% ID: SEL-D0076-OA-E01
% ESTADO: cotejada. Segunda pasada visual y borrador compilado revisados.
% Fuente: Selectividad/sel_2008_matematicas_aplicadas/Mat CCSS 4 - Examen y criterios de correción.pdf
% SHA-256: f1040adb7b9fb01c26219800e6212b9a26b93bb9d5e933743d18ee6ac26e163b
% Página: 1; opción A; ejercicio original 1.

\textbf{(3 puntos)} Un joyero fabrica dos modelos de anillos. El modelo A se hace con 1 gramo de oro y $1.5$ gramos de plata. El modelo B lleva $1.5$ gramos de oro y 1 gramo de plata. El joyero sólo dispone de 750 gramos de cada metal y piensa fabricar, al menos, 150 anillos del tipo B que ya tiene encargados. Sabiendo que el beneficio de un anillo del tipo A es de 50 \texteuro{} y del tipo B es de 70 \texteuro{}, \textquestiondown{}cuántos anillos ha de fabricar de cada tipo para obtener el beneficio máximo y cuál será éste?
